\documentclass[10pt,a4paper]{amsart}
\usepackage[margin=19mm]{geometry}
\usepackage{amsmath,amssymb,mathtools}
\usepackage[scr=rsfs,cal=euler]{mathalfa}
\usepackage{xfrac}
\usepackage[unicode,hidelinks,bookmarks=false]{hyperref}
\newcommand{\ie}{i.e.\ }
\newcommand{\cf}{cf.\ }
\newcommand{\resp}{respectively\ }
\newcommand{\Subsection}{\subsection}
\providecommand{\nobreakdash}{\nobreak\hbox{-}}
\setlength{\emergencystretch}{2em}
\multlinegap0pt
\usepackage{amssymb}
\usepackage{enumerate}
\usepackage{mathtools}
\newtheorem{lemma}{Lemma}
\newtheorem{proposition}{Proposition}
\newtheorem{claim}{Claim}
\title{Graphs whose Eulerian trails have unique labels}
\author{Donggyu Kim, Rose McCarty, Caleb McFarland}
\date{Source: arXiv:2603.02501v1 (CC BY 4.0)}
\begin{document}
\maketitle

\noindent\emph{Source and licence.} Graphs whose Eulerian trails have unique labels. Version arXiv:2603.02501v1 (CC BY 4.0), distributed by arXiv under the Creative Commons Attribution 4.0 International licence, http://creativecommons.org/licenses/by/4.0/, as recorded in the arXiv licence metadata for this exact version and shown on the abstract page https://arxiv.org/abs/2603.02501v1. Full licence notice: \texttt{licenses/arxiv-2603.02501-CC-BY-4.0.txt}.\par\smallskip

\noindent\textit{Editorial scope.} Counted unit: the article's proposition prop:3edgeConn (source0.tex lines 691--697). The article's complete original proof of it is delivered with it (source0.tex lines 698--746, one \texttt{proof} environment of 49 lines). No mathematics is written, repaired or re-derived: the statement and the proof are the article's own lines, in the article's own notation, and no retained line is substituted or re-flowed.  Retained uncounted as the source-local context the proof needs: lines 643--646.  Nothing was omitted from any retained span, so the package carries no omission record.  No retained line contains a citation, so the wrapper carries no bibliography and no bibliography entry.  No retained line contains a figure environment, an image-inclusion command or drawing code, so no image is delivered and the package declares no asset dependency.  Editorial formatting only, in addition: an AMS \texttt{amsart} preamble that restates the article's own declarations and macros used by the retained lines, namely the environments \texttt{lemma}, \texttt{proposition}, \texttt{claim} on separate counters, together with the standard packages those lines need.  The retained environments print in this document as Lemma 1, Proposition 1, Claim 1 in order of appearance, whereas the article prints the same items with the numbering of its own full document.  The complete native source of the article as distributed in the arXiv e-print for this exact version is bundled unchanged as \texttt{sources/195624/source0.tex}.
\begin{lemma}
\label{lem:splittingMinCut}
For any valid instance $(G, a, b)$ with at least one edge, there exists an edge $e$ with ends $a$ and $a'$ (possibly with $a=a'$) so that $(G_e, a', b)$ is a valid instance.
\end{lemma}

\begin{proposition}
\label{prop:3edgeConn}
Let $(G, a,b)$ be a valid instance, and let $\gamma$ be a group-labeling of $G$. Then exactly one of the following holds.\begin{enumerate}
    \item There exists a feasible circuit $L$ with $\gamma(L)\neq \gamma(L^{-1})$.
    \item There exists a shifting $\gamma'$ of $\gamma$ so that $\langle G, \gamma'\rangle\cong \mathbb{Z}_2^k$ for some $k\in \mathbb{N}$. Moreover, for any $e\in E(G)$, the set $\{\gamma'(L): L$ is a feasible circuit that contains $\vec{e}\}$ generates $\langle G, \gamma'\rangle$.
\end{enumerate}
\end{proposition}

\begin{proof}
    It is clear that both of the conditions cannot hold, so we just prove that at least one of the conditions holds. We prove this by induction on the number of edges in $G$. We may assume that $G$ has an edge as otherwise the second outcome trivially holds.

    By Lemma~\ref{lem:splittingMinCut}, there exists an edge $e \in E(G)$ with ends $a$ and $a'$ such that $(G_e, a', b)$ is a valid instance. If $a$ is incident to exactly two edges in $G-e$, then we write $e_1$ and $e_2$ for those two incident edges and $e_{1,2}$ for the new edge of $G_e$. We may assume that $\vec{e}_1$ is oriented towards $a$, $\vec{e}_2$ is oriented away from $a$, and $\vec{e}_{1,2}$ is oriented from the tail of $\vec{e}_1$ to the head of $\vec{e}_2$. Finally, let $\gamma_e$ denote the group-labeling of $G_e$ corresponding to $\gamma$. (First we forget about the label of $e$. Then, if $G_e \neq G-e$, we label $\vec{e}_{1,2}$ by $\gamma(\vec{e}_1)\gamma(\vec{e}_2)$.) 

    Now we apply induction to the valid instance $(G_e, a', b)$ and labeling $\gamma_e$. If the first outcome holds, then $(G, a,b)$ also has a feasible circuit $L$ so that $\gamma(L)$ has order greater than~$2$. Thus we may assume that there exists a shifting $\gamma_e'$ of $\gamma_e$ such that the second outcome holds for $(G_e, a', b)$. For convenience, we denote $\langle G_e, \gamma'_e\rangle$ by $\Gamma_e$. Let $k$ be the integer such that $\Gamma_e \cong \mathbb{Z}_2^k$. 
    
    We now define a corresponding shifting $\gamma'$ of $\gamma$ in~$G$. First perform the same shiftings (by group elements at vertices) as used to obtain $\gamma_e'$ from $\gamma_e$. Finally, if $G_e\neq G-e$, then perform one final shifting at $a$ in order to make $\vec{e}_1$ labeled by the identity.
    
    %Let $v_1, v_2, \ldots, v_n$ be vertices of $G_e$ and $\alpha_1, \alpha_2, \ldots, \alpha_n$ be group elements such that $\gamma_e'$ is obtained from $\gamma_e$ by first shifting at $v_1$ by $\alpha_1$, and then shifting at $v_2$ by $\alpha_2$, and so on. We now define a corresponding shifting $\gamma'$ of $\gamma$ in~$G$.

    Note that every orientation of an edge in $E(G-e)$, except for possibly $e_1$ and $e_2$ if they exist, has the same label according to $\gamma'$ and $\gamma_e'$. Moreover, if $e_1$ and $e_2$ do exist, then $\gamma'(\vec{e}_1)$ is the identity, and $\gamma'(\vec{e}_2)=\gamma_e'(\vec{e}_{1,2})$. Thus $\langle G - e, \gamma'\rangle = \Gamma_e$, which is isomorphic to $\mathbb{Z}_2^k$. Without loss of generality we assume that $\vec{e}$ is oriented away from $a$, and we write $\beta$ for $\gamma'(\vec{e})$. (If $e$ is a loop, then we orient $\vec{e}$ arbitrarily.) Note that $\langle G, \gamma'\rangle$ is the group generated by $\Gamma_e \cup \{\beta\}$.

    We first show that the set $\{\gamma'(L) : L$ is a feasible circuit that contains $\vec{e}\}$ generates $\langle G, \gamma'\rangle$. We write $\mathcal{L}$ for the collection of all feasible circuits of $(G, a, b)$ whose first arc is $\vec{e}$. 

    \begin{claim}
    \label{clm:ecircuit}
        The set $\{\gamma'(L): L \in \mathcal{L}\}$ generates $\langle G, \gamma'\rangle$.
    \end{claim}
    \begin{proof}
        Because $\Gamma_e \cup \{\beta\}$ generates $\langle G, \gamma'\rangle$, it suffices to show every element in $\Gamma_e \cup \{\beta\}$ is generated by $\{\gamma'(L): L \in \mathcal{L}\}$.
    
        If $e_1$ and $e_2$ exist, then we apply the inductive statement to the edge $f = e_{1,2}$. Otherwise we apply the inductive statement to an edge $f\neq e$ which is incident to $a$. (We may assume that such an edge exists since if $e$ is the only edge of $G$, then it is a loop at the vertex $a=b$, and the claim holds.) Thus, by the inductive statement, the group generated by $\{\gamma_e'(L): L$ is a feasible circuit of $(G_e, a', b)$ that contains $\vec{f}\}$ is $\Gamma_e$. In both cases, it follows that $\Gamma_e$ is generated by the $\gamma'$-labels of the circuits $L$ of $G$ such that there exist trails $T_1, T_2$ so that $T_1 L T_2$ is an Eulerian trail from $a$ to $b$ in $G$, the first edge of $T_1$ is $\vec{e}$, and $L$ contains at least one edge that is incident to $a$. 

        Consider such trails $T_1, L, T_2$ as above. Let $L_1$ and $L_2$ be subtrails of $L$ (possibly with one of $L_1, L_2$ being empty) such that $L=L_1 L_2$ and $a$ is the last vertex of $L_1$. Then both $T_1 L_1$ and $T_1 L_2^{-1}$ are in $\mathcal{L}$. Note that \begin{align*}
        (\gamma'(T_1 L_1))^{-1}\gamma'(T_1 L_2^{-1})&=\gamma'(L_1)^{-1}\gamma'(T_1)^{-1}\gamma'(T_1)\gamma'(L_2)^{-1}\\&=\gamma'(L_1)^{-1}\gamma'(L_2)^{-1}\\&=\gamma'(L_1)\gamma'(L_2),
        \end{align*}where the last line holds because $\gamma'(L_1)$ and $\gamma'(L_2)$ have order at most~$2$. (To see this, recall that $\vec{e}$ is in $T_1$ and the other arcs generate $\Gamma_e$, which is isomorphic to $\mathbb{Z}_2^k$.) Since $\gamma'(L_1)\gamma'(L_2)= \gamma'(L)$, it follows that the $\gamma'$-labels of the circuits in $\mathcal{L}$ generate $\Gamma_e$.
    
        Moreover, for any $L \in \mathcal{L}$, there exists $\alpha \in \Gamma_e$ so that $\gamma'(L)=\beta\alpha$. Since we have shown that $\alpha$ is generated by the $\gamma'$-labels of the circuits in $\mathcal{L}$, it follows that $\beta$ is generated by them as well.
    \end{proof}

    Since every circuit in $\mathcal{L}$ is feasible for $(G, a, b)$, Claim~\ref{clm:ecircuit} implies that $\{\gamma'(L): L$ is a feasible circuit of $(G, a, b)$ that contains $\vec{e}\}$ generates $\langle G, \gamma'\rangle$.

    We now show that $\langle G, \gamma'\rangle$ is isomorphic to $\mathbb{Z}_2^{k'}$ for some $k' \in \mathbb{N}$. Clearly this holds if $\beta \in \Gamma_e$, so assume otherwise. Now consider the $\gamma'$-label of a circuit $L \in \mathcal{L}$. It must be of the form $\beta\alpha$ for some $\alpha \in \Gamma_e$. If any of these elements $\beta\alpha$ has order greater than~$2$, then the first outcome of the lemma holds, and we are done. So we may assume that each of these elements has order at most~$2$. Then any two of these elements $\beta\alpha$ and $\beta\alpha'$ commute since
        \begin{align*}
        (\beta\alpha)(\beta\alpha') &= (\alpha\beta^{-1})(\beta\alpha')
        = \alpha\alpha'
        =\alpha'\alpha
        = (\alpha'\beta^{-1})(\beta\alpha)
        = (\beta\alpha')(\beta\alpha).
        \end{align*}
    Any group which is generated by a finite number of commuting elements of order $2$ is isomorphic to $\mathbb{Z}_2^{k'}$ for some $k' \in \mathbb{N}$; thus it follows that $\langle G, \gamma'\rangle$ is isomorphic to~$\mathbb{Z}_2^{k+1}$.

    We now show that for any edge $f \in E(G - e)$, the set $\{\gamma'(L) : L$ is a feasible circuit of $(G,a,b)$ that contains $\vec{f}\}$ generates $\langle G,\gamma'\rangle$. Fix $f \in E(G - e)$, and let $\Gamma^f$ denote the group generated by the labels of feasible circuits of $(G,a,b)$ containing $\vec{f}$. Because each feasible circuit of $(G_e, a', b)$ ``lifts'' to a feasible circuit of $(G, a, b)$ of the same label, $\Gamma^f \supseteq \Gamma_e$. Thus we may assume that $\beta \not\in \Gamma_e$, and it suffices to show that $\beta \in \Gamma^f$.

    It then suffices to show that there exists a feasible circuit which contains both $e$ and $f$ as underlying undirected edges. Since the group is isomorphic to $\mathbb{Z}_2^{k+1}$, the order and orientation do not matter. Once we show that such a feasible circuit exists, it follows that $\beta$ is also generated by the set $\{\gamma'(L): L$ is a feasible circuit that contains $\vec{f}\}$, as we already know that the labels of all other edges besides $e$ are generated.
    
    Now we argue that such a feasible circuit exists. Since $G$ is $2$-edge-connected, every pair of edges of $G$ is in a common circuit by Menger's Theorem for edge-disjoint paths. Let $L$ be an edge-maximal circuit of $G$ which contains both $e$ and $f$. We claim that $L$ is feasible. By edge-maximality, every component of $G-E(L)$ which contains an edge has a vertex of odd degree. Thus, if $a=b$, then $L$ is an Eulerian circuit and we are done. If $a \neq b$, then $a$ and $b$ are in the same component of $G-E(L)$, and every other component is just an isolated vertex. Let $T$ be an Eulerian trail from $a$ to $b$ within the non-trivial component. Let $v$ be a vertex along both $T$ and $L$, and let $T = T_1T_2$ such that the head of $T_1$ equals the tail of $T_2$ equals $v$ (possibly with one of $T_1,T_2$ being empty). Then $T_1LT_2$ is an Eulerian trail of $G$, thus showing that $L$ is feasible, as desired. This completes the proof of Proposition~\ref{prop:3edgeConn}.
\end{proof}

\end{document}
